\documentclass[11pt]{article}
\usepackage[margin=.9in]{geometry}
\usepackage{amsmath,amssymb,amsthm}
\usepackage[T1]{fontenc}
\usepackage{lmodern,microtype,hyperref}
\newtheorem{theorem}{Theorem}
\newtheorem{lemma}[theorem]{Lemma}
\newtheorem{proposition}[theorem]{Proposition}
\newtheorem{corollary}[theorem]{Corollary}
\theoremstyle{remark}\newtheorem{remark}[theorem]{Remark}
\newcommand{\norm}[1]{\left\lVert#1\right\rVert}
\DeclareMathOperator{\diag}{diag}\DeclareMathOperator{\ran}{ran}
\newcommand{\R}{\mathbb R}\newcommand{\C}{\mathbb C}
\begin{document}
\section{Strict support localization of the homogeneous classical action}
\label{cl-section}

Fix odd $L$, dimension four, $M=L^j$, $\eta=M^{-1}$ and
$0<\alpha<1$. Work first on compatible finite periodic regulators.
Constants below are independent of $M$ and volume. They may depend on
$L$, color, $\alpha$, the requested positive spatial exponent and an
arbitrarily prescribed finite support-cardinality reserve. The product
background radius may shrink with these fixed choices. Fine transposes
retain the physical measure $\eta^4$.

The concrete cell-vertex estimates preceding this section supply the
exact simultaneous equations
\begin{align}
 A&=HB-HR(A)-\mathcal C S^T f_\eta(SA)
                         +\mathcal C V(A,m),\notag\\
 m&=qB-qR(A)+(SH)^T f_\eta(SA)-H^T V(A,m),
 &V(A,m)&=DR(A)^Tm.                                      \label{cl-fixed-point}
\end{align}
Here $\mathcal F(A)=QA+R(A)$ is the full corrected fixed-reference
constraint, $H$ is the matched flat minimizing extension,
$\mathcal C$ the flat constrained inverse, $q=q_M$, and $S$ the local
first-difference jet. All three nonlinear primitives start at total
degree two. Rearranging the second equation gives the multiplier inverse
of the constrained stationary equation; the first then gives its exact
constraint chart. Thus \eqref{cl-fixed-point} is equivalent to the
previously proved small classical branch. It uses the physical adjoint
of the corrected constraint, not the nonuniform adjoint of the full
ordinary rooted average. The scalar action equality supplied by the
ordinary-constraint gauge transfer will be used after solving these
equations.

\subsection{Geometry and the support-family norm}
\label{cl-geometry-section}

Let $d$ be current-cell graph distance. For a nonempty finite set $X$,
let $\tau(X)$ be the minimum number of edges in a connected graph
containing $X$, allowing unmarked intermediate vertices. This graph
Steiner length differs from the metric minimum spanning-tree length
$\ell$ used in the Gaussian expansion, but
\begin{equation}
 \tau(X)\le\ell(X)\le2\tau(X),\qquad |X|\le\tau(X)+1.       \label{cl-metric-comparison}
\end{equation}
The second upper bound follows by traversing a connected minimizing
tree and shortcutting its tour in the graph metric. The crucial rooted
inequalities are
\begin{align}
 \tau(\{i\}\cup X)&\le d(i,j)+\tau(\{j\}\cup X),\notag\\
 \tau\left(\{i\}\cup\bigcup_{r=1}^nX_r\right)
 &\le\tau(\{i,j_1,\ldots,j_n\})
                         +\sum_{r=1}^n\tau(\{j_r\}\cup X_r).
                                                               \label{cl-root-gluing}
\end{align}
The auxiliary roots $j_r$ need not be final support labels. This is why
the support algebra uses $\tau$: the analogous inequality for metric
spanning trees with deleted auxiliary vertices need not hold. All these
statements remain valid on periodic graphs.

Use the smooth fixed-physical-scale partition of unity and the complete
cell spaces from the concrete vertex construction. Their field types
are $\mathcal N_\alpha$, $C^\alpha$, bounded fine fields, and
fixed-dimensional current fields. Localization $P_i$ and reconstruction
$\mathcal R$ are uniformly bounded, and $\mathcal R P=I$.
In particular, input pieces to $\mathcal C S^T$ remain complete
$C^\alpha$ fields; no sharp cut is differentiated. For each relevant
linear propagator $T$, the proved weighted complete-field estimate gives
\begin{equation}
 \|T_{ij}\|\le C_Te^{-\kappa_0d(i,j)},\qquad
 L_T(\kappa):=\sup_i\sum_j e^{\kappa d(i,j)}\|T_{ij}\|<\infty
 \quad(0<\kappa<\kappa_0).                                  \label{cl-linear-blocks}
\end{equation}
The diagonal block of $\mathcal C S^T$ is controlled by the proved
equal-exponent Schauder estimate, rather than an absolute singular-kernel
integral. These block estimates are uniform in refinement and volume.

On the original current-background product ball of radius $b$, a support
family is a collection $u_{i;X}(B|_X)$ of bounded holomorphic functions
with values in the cell space at $i$. The support $X$ is nonempty and
records original background cells. For $a\ge0$ put
\begin{equation}
 \|u\|_{\kappa,a;b}
 :=\sup_i\sum_X e^{\kappa\tau(\{i\}\cup X)+a|X|}
       \sup_{\|B|_X\|_\infty<b}\|u_{i;X}(B|_X)\|.             \label{cl-family-norm}
\end{equation}
For several field types, take their maximum in each coefficient norm
\emph{before} the support sum: the joint family norm is
$\sup_i\sum_X e^{\kappa\tau(\{i\}\cup X)+a|X|}
\max_t\|u^t_{i;X}\|$, with the displayed background supremum understood.
It is not the maximum of the separate type-family norms. Taylor
operator norms below use this same joint cell space with its maximum
norm. A bound for a joint linear source may be taken as the sum of its
component block bounds. These are fixed finite sums.
This is a Banach space. Terms with the same support are added in that
slot; evaluation $\sum_Xu_{i;X}(B|_X)$ converges absolutely. All
coefficients remain functions of the original variables. No nonlocal
background-dependent minimizer is hidden in their definition.

\begin{lemma}[Support substitution]\label{cl-substitution}
A linear propagator obeying \eqref{cl-linear-blocks} has norm at most
$L_T(\kappa)$ in \eqref{cl-family-norm}. Suppose a nonlinear cell map has
Taylor coefficients, including $1/n!$, whose mixed-insertion norms obey
\begin{align}
 v_n&:=\sup_i\sum_{j_1,\ldots,j_n}
 e^{\kappa\tau(\{i,j_1,\ldots,j_n\})}
       \|N_{n;i;j_1,\ldots,j_n}\|,\notag\\
 \Phi(t)&:=\sum_{n\ge2}v_nt^n<\infty\quad(0\le t<t_0).
                                                               \label{cl-vertex-majorant}
\end{align}
Assign a substituted term the union of its input supports. Then
\begin{align}
 \|N(u)\|_{\kappa,a;b}&\le\Phi(\|u\|_{\kappa,a;b}),\notag\\
 \|N(u)-N(v)\|_{\kappa,a;b}&\le\Phi'(r)\|u-v\|_{\kappa,a;b}
 \quad(\|u\|,\|v\|\le r<t_0).                              \label{cl-substitution-bound}
\end{align}
The majorant is independent of the prescribed cardinality reserve $a$.
\end{lemma}

\begin{proof}
For the linear assertion attach a path from the input root $j$ to the
output root $i$ and use the first inequality in
\eqref{cl-root-gluing}. For a nonlinear term use the second inequality
and
$|\bigcup_rX_r|\le\sum_r|X_r|$.
The absolute sum over all supports and ordered input-cell slots then
factorizes into the product of their family norms, leaving $v_n$.
Summing over $n$ proves the first bound. Telescoping an $n$-linear
expression gives the factor $n r^{n-1}$ and proves the second.
Repeated supports, repeated labels and overlapping supports are all
included in these sums. A Taylor coefficient at the fixed flat
background introduces no additional variable support, so it adds no
cardinality cost.
\end{proof}

The concrete finite-range analytic vertices have
\begin{equation}
 v_n\le M_1
       \left(\frac{eC_PD_Re^{\kappa R}}{r_0}\right)^n.        \label{cl-local-cauchy}
\end{equation}
Here $R$ is a uniform current-cell dependency radius, $D_R$ the fixed
number of admissible input cells in one slot, $C_P$ the reconstruction
norm, and $r_0$ a uniform analytic radius. Indeed multivariable Cauchy
with radius $r_0/(nC_P)$ bounds the derivative by
$M_1(nC_P/r_0)^n$; division by $n!$ and $n^n/n!\le e^n$ bounds its
Taylor coefficient. There are at most $D_R^n$ nonzero ordered cell
tuples, with a star of length at most $nR$ through the output.
This counts current cells, not $M^4$ independent fine coordinates.
For $V(A,m)$ it uses the previously proved holomorphy in the physical
absolute-row norm of $DR(A)^T$; ordinary bounded-field holomorphy of
$R$ alone would not imply that input.

\subsection{A convergent support-valued stationary solution}
\label{cl-contraction-section}

\begin{theorem}[Uniform strict-support fixed point]
\label{cl-fixed-point-theorem}
For every fixed finite $a\ge0$ and sufficiently small positive
$\kappa$, there is $b_*>0$, independent of $M$ and volume, such that
the solution of \eqref{cl-fixed-point} has support families satisfying
\begin{align}
 \|(A,m)\|_{\kappa,a;b}&\le C_a b,\notag\\
 \|(A-HB,m-qB)\|_{\kappa,a;b}&\le C_a b^2
 \qquad(0<b\le b_*).                                      \label{cl-solution-bound}
\end{align}
Every coefficient of the second line has Taylor degree at least two.
\end{theorem}

\begin{proof}
Apply the preceding lemma to the three nonlinear primitives in
\eqref{cl-fixed-point} and then their linear propagators. Their combined
positive coefficient majorant can be bounded coefficientwise by
\begin{equation}
 \Phi(t)\le\frac{v(ct)^2}{1-ct},\qquad
 \Phi'(t)\le8vc^2t\quad(ct\le1/2),                          \label{cl-majorant-constants}
\end{equation}
after increasing fixed $v,c$. The singleton source
$B_{i;\{i\}}=B_i$ has norm $e^ab$. If $L_0$ bounds the linear map
$B\mapsto(HB,qB)$, choose
\begin{equation}
 r_*\le\min\{(2c)^{-1},(16vc^2)^{-1}\},\qquad
 b_*\le\frac{r_*}{2L_0e^a}.                                \label{cl-radius}
\end{equation}
Also shrink $r_*$ so that evaluation and reconstruction stay in the
previously established classical branch's uniqueness ball; their norms
are uniformly bounded. The support map is a contraction with constant
at most $1/2$ on the ball of radius $2L_0e^ab$. This proves the first
estimate. Applying the quadratic majorant to its solution proves the
second. Degree induction in Picard iteration proves the vanishing
orders coefficientwise.

Evaluation is continuous and respects every absolutely convergent
substitution. The evaluated solution therefore solves the exact
equations \eqref{cl-fixed-point}. The compatible cell image is preserved
by their right-hand side, and uniqueness in the common small domain
identifies the result with the proved classical minimizer and
multiplier. Picard iteration also defines an absolutely convergent
rooted-tree expansion: local Taylor coefficients are vertices, flat
cell propagators are edges, and the leaves carry the original $B_i$.
All depths, orders and repeated supports are already summed by the
Banach contraction. No additional factor counting the number of RG
levels occurs.
\end{proof}

\subsection{Scalar action and symmetric support Hessians}
\label{cl-action-section}

The established envelope identity is
\begin{equation}
 DS_M(B)=m(B),\qquad S_M(0)=0,\qquad D^2S_M(0)=q_M.          \label{cl-envelope-input}
\end{equation}
The exact gauge transfer identifies this scalar with the ordinary
homogeneous constrained action. Write $n=m-qB$. On a finite regulator,
\begin{equation}
 S_M(B)-\tfrac12\langle B,qB\rangle
    =\int_0^1\sum_i\langle B_i,n_i(tB)\rangle\,dt.           \label{cl-envelope}
\end{equation}
The current pairing has unit cell volume and fixed fiber dimension.
Define strictly supported scalar coefficients by
\begin{equation}
 W_Y(B|_Y)=
 \sum_{\{i\}\cup X=Y}\int_0^1
       \langle B_i,n_{i;X}(tB|_X)\rangle\,dt.                \label{cl-action-coefficients}
\end{equation}
They are holomorphic and begin at order three.

\begin{proposition}[Strict scalar action localization]
\label{cl-action-theorem}
For $0<\mu<\kappa$ and every prescribed finite $a\ge0$, on a fixed
small product ball,
\begin{align}
 S_M(B)&=\tfrac12\langle B,q_MB\rangle+\sum_{Y\ne\varnothing}W_Y(B|_Y),
                                                               \label{cl-action-expansion}\\
 \sup_b\sum_{Y\ni b}e^{\mu\tau(Y)+a|Y|}\|W_Y\|_r
       &\le C_{\mu,a}r^3.                                    \label{cl-action-norm}
\end{align}
Here $\|W_Y\|_r$ is the supremum on the complex product ball of radius
$r$. Each Hessian $B_Y=D^2W_Y$, extended by zero outside cell rows and
columns in $Y$, is individually complex symmetric and depends only on
$B|_Y$. With a fixed cardinality spare $\epsilon>0$,
\begin{equation}
 \sup_i\sum_{Y\ni i}\sum_{j\in Y}
 e^{\mu\tau(Y)+(a-\epsilon)|Y|}
            \|(B_Y)_{ij}\|_r\le C r.                       \label{cl-hessian-support}
\end{equation}
\end{proposition}

\begin{proof}
For $b\in Y=\{i\}\cup X$, one has $d(i,b)\le\tau(Y)$ and
$|Y|\le|X|+1$. Thus
\begin{equation}
 e^{\mu\tau(Y)+a|Y|}
 \le e^a e^{\kappa\tau(Y)+a|X|}
                   e^{-(\kappa-\mu)d(i,b)}.                \label{cl-reroot}
\end{equation}
Summing first over $X$ and then $i$, and using
$\sup_b\sum_i e^{-\delta d(i,b)}\le\sum_{x\in\mathbb Z^4}e^{-\delta|x|}$,
proves \eqref{cl-action-norm} from \eqref{cl-solution-bound} and
\eqref{cl-action-coefficients}. This is an anchored sum over all supports
containing an arbitrary cell, not only the original output root.

Use the action estimate on radius $2r$ and take derivatives on radius
$r$. Cauchy and polarization give each cell block the dimension-free
bound
\begin{equation}
 \|(D_iD_jW_Y)\|_r\le\frac{2e^2}{r^2}\|W_Y\|_{2r}.         \label{cl-hessian-cauchy}
\end{equation}
The row sum contributes at most $|Y|$, absorbed by
$\sup_{n\ge1}ne^{-\epsilon n}<\infty$ in the reserve. This proves
\eqref{cl-hessian-support} and justifies termwise differentiation.
Scalar holomorphy gives symmetry of each coefficient. Rows, columns and
background support are all in the same $Y$ by construction.
\end{proof}

Taking $\mu\ge2(\kappa_{\rm G}+2m_{\rm G})$ converts
\eqref{cl-hessian-support} into the strict precision-support norm with
metric-tree weight $e^{(\kappa_{\rm G}+2m_{\rm G})\ell(Y)}$ used by the
Gaussian density algebra. The Gaussian exponents can be chosen positive
and sufficiently small compared with the available flat decay. The
smallness of the precision norm is then obtained by shrinking the
background radius, not by an unsupported claim about a large original
source domain.

\subsection{The actual current pivot chart and retained precision}
\label{cl-chart-section}

Use the fixed-$L$ ordinary pivot and coordinate maps already identified
for the physical current step. With $a=\log V/i$, set
\begin{equation}
 F_{\rm cur}(a,z)=\frac1i\log\!\left(e^{i\Phi_V(z)}e^{ia}\right),
 \qquad F_{\rm cur}(a,0)=a.                                  \label{cl-current-chart}
\end{equation}
They are holomorphic on a fixed joint product neighborhood, have a fixed
current-cell dependency radius $R_c$, and satisfy
\begin{equation}
 \|F_{\rm cur}(a,z)\|_\infty
             \le C_L(\|a\|_\infty+\|z\|_\infty).            \label{cl-current-chart-bound}
\end{equation}
Let $G_{\rm cur}(a,z)$ be the actual local exponential gauge action on
$e^{i\Phi_V(z)}V$, not a replacement relative quadratic penalty. Its
full analytic expression is finite range at fixed $L$. At the physical
rooted-contour background its quadratic form is the matched gauge term;
away from that locus its actual local curvature terms remain included.
Put
\begin{equation}
 T_M(a,z)=S_M(F_{\rm cur}(a,z))+G_{\rm cur}(a,z).             \label{cl-pulled-action}
\end{equation}

The scalar family for $S_M$ includes its quadratic baseline: split
$\frac12\langle B,qB\rangle$ into pair-supported scalar coefficients.
Exponential locality of $q$ gives their anchored support norm of order
$r^2$ with any fixed cardinality reserve and a sufficiently small
positive spatial exponent. Adding \eqref{cl-action-coefficients}
therefore gives a scalar support family for the entire $S_M$.

For a support $Y$ let $\Gamma Y$ be the union of the radius-$R_c$
dependency neighborhoods of its cells. Fixed constants $N_c,c_c$ obey
\begin{equation}
 |\Gamma Y|\le N_c|Y|,\qquad
 \tau(\Gamma Y)\le\tau(Y)+c_c|Y|.                           \label{cl-support-enlargement}
\end{equation}
Attach a fixed finite neighborhood tree at each point of a minimizing
connected graph for $Y$ to prove the second inequality. Periodic
identifications only reduce this count. Hence
\begin{equation}
 e^{\sigma\tau(\Gamma Y)+a_{\rm out}|\Gamma Y|}
 \le e^{\sigma\tau(Y)+(\sigma c_c+a_{\rm out}N_c)|Y|}.       \label{cl-composition-weight}
\end{equation}
Choose the input cardinality reserve to pay this fixed cost, and shrink
the joint radius so that \eqref{cl-current-chart-bound} remains in the
source-action ball. Substitute the chart in each scalar coefficient and
assign support $X=\Gamma Y$. A root in $\Gamma Y$ can be moved to a
point of $Y$ in at most a fixed number of ways. Adding the local gauge
coefficients gives the exact expansion
\begin{equation}
 T_M(a,z)=\sum_X T_X(a|_X,z|_X),                             \label{cl-pulled-support-expansion}
\end{equation}
with a uniform anchored scalar norm on the fixed joint product domain
and any prescribed fixed output cardinality reserve. Both the original
background and retained-coordinate labels remain explicit in $X$.

\begin{corollary}[Strict support of the physical retained precision]
\label{cl-precision-theorem}
Define
\begin{equation}
 A_X(a|_X)=D_z^2T_X(a|_X,0)-D_z^2T_X(0,0).                   \label{cl-precision-coefficients}
\end{equation}
These are individually complex-symmetric retained matrices with rows,
columns and background dependence in $X$. For some $\sigma>0$ and all
sufficiently small $b$,
\begin{align}
 A_M(V)&=A_M(1)+\sum_X A_X((\log V/i)|_X),                  \label{cl-precision-expansion}\\
 \sup_i\sum_{X\ni i}\sum_{j\in X}
 e^{\sigma\tau(X)}\|(A_X)_{ij}\|_b&\le C_L b.             \label{cl-precision-norm}
\end{align}
The baseline is exactly the matched flat retained precision
$A_M(1)=K_L^T(q_M+T_L^TT_L)K_L$.
\end{corollary}

\begin{proof}
Apply Cauchy with fixed margins to the joint domain of
\eqref{cl-pulled-support-expansion}, twice in $z$ and once in $a$.
Integrating the last derivative along $ta$ gives $O(b)$ in each
coefficient. Its row sum costs $|X|$, paid by an available cardinality
reserve as in \eqref{cl-hessian-support}. The baseline sum is the flat
Hessian of the exact pulled action and hence the identified retained
precision. The matrices are Hessians of scalar coefficients, which
proves their individual symmetry and supports.
\end{proof}

This argument retains both nonlinear constraint curvature and
exponential-coordinate curvature. It does not identify a freely chosen
independent source current with the physical Hessian. The full product
background extension is an analytic Hessian family; only the physical
background is asserted to be stationary. Choosing
$2(\kappa_{\rm G}+2m_{\rm G})<\sigma$ and then shrinking $b$ makes the
strict norm as small as the stated Gaussian support algebra requires.
The baseline covariance is background independent and has the uniform
spectrum and decay proved in the flat theory.

\subsection{Localized cubic interaction and Gaussian-cell regrouping}
\label{cl-interaction-section}

For $\|z\|_\infty\le R$ and $gR$ inside a fixed analytic margin define
\begin{align}
 U_X(a;z)=-g^{-2}\bigl[&T_X(a,gz)-T_X(a,0)
          -gD_zT_X(a,0)[z]\notag\\
          &-\tfrac12g^2D_z^2T_X(a,0)[z,z]\bigr].             \label{cl-cubic-coefficients}
\end{align}
The sign corresponds to $\exp(-g^{-2}T_M)$.
Taylor's integral remainder, or one-variable Cauchy along $z/R$, yields
\begin{equation}
 \sup_i\sum_{X\ni i}e^{\sigma\tau(X)+a_{\rm out}|X|}
 \sup_{\|a\|<b,\,\|z\|\le R}|U_X(a;z)|\le CgR^3.          \label{cl-cubic-norm}
\end{equation}
No factor $|X|^3$ is needed: the full direction $z|_X$ has product norm
at most $R$. At a physical stationary background the sum of the linear
terms in \eqref{cl-cubic-coefficients} vanishes. This cancellation is
not asserted for every point of the analytic product-background domain.

For the fixed retained product cutoff set
$\chi_X=\prod_{i\in X}\chi_i$ and
$w_X=\chi_X(e^{U_X}-1)$. Then
$v_X:=\sup|w_X|\le e^{\|U_X\|}\|U_X\|$.
Take $a_{\rm out}>4+\log\widehat K$ with a positive spare, where
$\widehat K$ is the chosen fixed Gaussian density factor, and choose
$2\kappa_{\rm G}\le\sigma$. Equations
\eqref{cl-metric-comparison} and \eqref{cl-cubic-norm} give
\begin{equation}
 \sup_i\sum_{X\ni i}v_X\widehat K^{|X|}
                   e^{4|X|+\kappa_{\rm G}\ell(X)}
        \le C e^{CgR^3}gR^3.                                \label{cl-polymer-input}
\end{equation}
Thus this is the required localized interaction norm, not only an
extensive unweighted cubic remainder. A logarithmically growing $R$ as
$g\downarrow0$ is compatible with both the analytic restriction on $gR$
and smallness of $gR^3$. The separate Gaussian bad-block tail and the
explicit density-algebra smallness conditions must still be imposed.
No all-bond nonlinear sharp cutoff is identified with the retained
product window by this argument.

The support algebra above uses unit current cells. The Gaussian theorem
may use one cell per fixed $L$-block, with
$D=d(L^d-1)\dim\mathfrak g$ retained real coordinates in the usual pivot
convention. Let $\pi$ map a unit current cell to its $L$-block and regroup
each support $Y$ to $\pi(Y)$. Then
\begin{equation}
 \tau_{\rm block}(\pi(Y))\le\tau_{\rm unit}(Y),\qquad
 |\pi(Y)|\le|Y|.                                           \label{cl-block-regrouping}
\end{equation}
Every unit edge either stays in one block or crosses to an adjacent
block. A block-anchored support sum is bounded by the sum of the unit
anchored sums over its at most $L^d$ cells. For matrices, bound a block
matrix norm by the sum of its constituent unit-cell blocks, then sum the
at most $L^d$ row cells. Thus scalar and absolute-row support norms
change by only fixed $L$-dependent factors. Maximum-cell and Euclidean
block norms likewise differ by fixed factors; reduce a background or
cutoff radius accordingly. No dimension depending on $M$ is introduced.
The Gaussian density constants in \eqref{cl-polymer-input} can therefore
use the fixed retained dimension $D$ after this explicit regrouping.

\subsection{Scope of the classical input}
\label{cl-scope}

The construction expands about a fixed flat baseline in the original
current variables. The arbitrary finite cardinality reserve is paid by
a smaller uniform product radius; spatial decay alone would not pay
the Gaussian factor $\widehat K^{|X|}e^{4|X|}$. Fine dimension is removed
by complete cell norms, the physical measure, and the actual
physical-adjoint row bound. The scalar ordinary action is identified by
the proved gauge transfer, while its current pivot chart has fixed
ratio $L$.

The conclusions supply strict-support precision and classical
interaction inputs for the homogeneous classical step at every depth
$M$, on the indicated small branch. They do not establish closure of
inherited quantum polymer terms, arbitrary independent-current or
mixed-history domains, a running-coupling theorem, or an interacting
continuum limit. Finite compatible periodic regulators are primary;
the graph estimates are uniform, but source chart incidence still uses
the admissible current tori assumed in its proof. No thermodynamic
consistency or separate gauge invariance of an arbitrarily regrouped
individual coefficient is asserted.

\end{document}
